\section{Finite-Enclosure Calculation Details}
\label{app:finite-calculations}

This appendix contains the scalar and selector calculations suppressed from
the reader-facing Method~3 chapter.  The body retains every theorem interface
and every geometric implication.

\subsection{Four direct high-radial outputs}
\label{subsec:calc-four-direct-outputs}

\begin{proof}[Proof of Proposition~\ref{prop:new-four-direct-outputs}]
Since $c>1/2$, any containing V triangle contains the radial midpoint.
Therefore $a+b\le1$, and the supercritical cell of Proposition
\ref{prop:new-exact-local-set} is absent.  On the cap $b=1-a$, the $T$-cell
inequality becomes
\[
 M(c-M)\le0,
\]
where $M=\max\{a,b\}$.  Thus the linear cap is feasible exactly when
$a\le1-c$ or $a\ge c$.

Off the linear cap, a maximal point lies on $F_L=0$ or $F_T=0$.  The selected
$L$-cell equation has the two ordered roots $\ell(c),r(c)$; the selector
$q\le0$ retains the smaller coordinate.  The ordered halves of $F_T=0$ give
$Q_+$ and $Q_-$.  Their equality boundary $a=b$ is $h(c)$ in the low-radial
range, while their $q=0$ boundaries are
$(\ell(c),r(c))$ and $(r(c),\ell(c))$ in the high range.  These conditions
produce (6.66)--(6.67).

The formal other root in the $Q_+$ quadratic violates the geometric selector
$c\le2\max\{a,b\}$ and is discarded.  At the high-range junction, the
selected $T$-cell point has value $r(c)$, whereas the adjacent selected
$L$-cell point has value $\ell(c)$; this is the genuine downward jump stated
above.
\end{proof}

\subsubsection{Low-root and threshold calculations}

Let $c=1-e$ with $0<e<1-\sqrt3/2$.  The selected low root
$\epsilon(e)=\ell(1-e)$ satisfies the quadratic
\[
 p_e(z)=z^2-(1-e)z+(1-e)^2-(1-e)^4=0.
 \tag{6.68}
\]
Substitution gives
\[
 p_e(e)=e(1-3e+4e^2-e^3)>0.
\]
Since $e<(1-e)/2$, the point $e$ lies below the smaller root.  For
$U=2e/(1-2e)$,
\[
 p_e(U)=
 -\frac{e^2}{(1-2e)^2}
 (4e^4-20e^3+37e^2-28e+3).
\]
The polynomial in parentheses is positive on the required interval, proving
\[
 e<\epsilon(e)<\frac{2e}{1-2e}.
\]
For $e\le1/8$,
\[
 p_e(2e+5e^2)=e^2(3e+4)(8e-1)\le0,
\]
which proves (6.6).

Now let $A>\epsilon(e)$ and $C\ge1-e$.  The linear branch of
Proposition~\ref{prop:new-four-direct-outputs} is no longer available.  On
the selected constant and $Q_-$ pieces, the outgoing reach is at most
$\epsilon(e)$.  On the $Q_+$ piece, the selected quadratic has its outgoing
coordinate below the same root.  Thus $B\le\epsilon(e)$, proving the direct
threshold Lemma~\ref{lem:new-direct-threshold} with no following-edge map.

\subsubsection{The strict-supercritical envelope}
\label{subsec:calc-strict-supercritical}

Fix $0\le c<1/2$ and maximize the outgoing boundary reach over strict
supercritical pairs.  The maximum lies on the $S$-cell frontier and at the
boundary where the incoming coordinate tends to the complementary value.
Eliminating that coordinate from $F_S=0$ gives
\[
 b^2-cb+2c-1=0.
\]
The larger selected root is
\[
 M_c^{\rm sup}=\frac{c+\sqrt{c^2-8c+4}}2.
\]
The corresponding equality state lies on the boundary $a+b=1$, so it is not
strictly supercritical.  Therefore every actual strict-supercritical role has
\[
 B<M_c^{\rm sup}.
\]
This strict nonattainment is the only supercritical output used in the T3-like
and Vd1 rescuer proofs.



\subsection{Quarter radial envelope}
\label{subsec:calc-quarter-envelope}

\begin{proof}[Proof of Lemma~\ref{lem:new-quarter-radial-envelope}]
Only the $L$ and $T$ cells of Proposition~\ref{prop:new-exact-local-set} can
occur.  In the $L$ cell, the selected radial root is the unique root of
\[
 P_{h_0}(c)=c^4-c^2+h_0c-h_0^2.
\]
At $c_0=1-h_0/4$,
\[
 P_{h_0}(c_0)
 =\frac{h_0}{256}
 (h_0^3-16h_0^2-240h_0+128)>0
\]
for $0<h_0\le1/2$.  Since
\[
 P_{h_0}(\sqrt3/2)=-(h_0-\sqrt3/4)^2\le0
\]
and $P_{h_0}$ is increasing on the selected root interval, the selected root
is below $c_0$.

In the $T$ cell put $s=p+h_0$ and $r=\sqrt{4s^2-3}$.  Write
\[
 \gamma=\frac{p-h_0}{s}=wr,
 \qquad0\le w<1.
\]
Then
\[
 p=\frac{s(1+wr)}2,
 \qquad h_0=\frac{s(1-wr)}2,
\]
and
\[
 c_T=\frac{s(1+wr)}{1+r}.
\]
For fixed $s,r$, the function $c_T+h_0/4$ increases in $w$, so
\[
 c_T+\frac{h_0}{4}
 <\frac{s(9-r)}8.
\]
Using $4s^2=r^2+3$,
\[
256-(r^2+3)(9-r)^2
=(1-r)(r^3-17r^2+67r+13)\ge0.
\]
Hence $c_T+h_0/4<1$.  The two cells are exhausted.
\end{proof}

\subsubsection{Vd corner margins}

In the exact Vd1/Vd2 corner form, let the boundary reaches be $a,b$, let
$t>0$, and put $d=\sqrt{t^2+t+1}$.  The own-radial reach is
\[
 c_0=\frac{d-a-tb}{t+1}.
\]
If the role has positive support on the forward adjacent arm, its far endpoint
there is
\[
 u_+=\frac{d-a-tb-1}{t}.
\]
Since $d<t+1$,
\[
 c_0<1-\frac{a+tb}{t+1}
 \le1-\min\{a,b\},
 \tag{6.70}
\]
and
\[
 u_+<1-b-\frac at<1-b.
 \tag{6.71}
\]
Thus lower boundary demands $a\ge A$, $b\ge H$ imply
\[
 c_0<1-\min\{A,H\},
 \qquad
 u_+<1-H.
\]
Reflection supplies the backward adjacent-arm statements.  These two
one-line inequalities are exactly the radial separations used in the one-Vd
placement proof.


\subsection{CE1 reverse-path certificate}
\label{subsec:calc-ce1-reverse-path}

\begin{proof}[Proof of Proposition~\ref{prop:new-ce1-direct-certificate}]
The sign difference in (6.12) gives
\[
 RD>wA,
\]
so the center exit on $r_3$ is $m=A/R$.  Moreover
\[
 0<A<P<\frac18,
 \qquad 0<D<\frac R2,
 \qquad 0<m<\frac w2<\frac12.
 \tag{6.15}
\]
We first record
\[
 \epsilon(A)<X,
 \qquad h_0>A.
 \tag{6.16}
\]
Indeed, $A+wD<P$ and $D=R-X$ give $A<wX-\eta$.  Convexity of
\[
 wX-\frac{X}{2(1+X)}
\]
on $R/2<X<R$ gives
\[
 A<\frac{X}{2(1+X)}.
\]
Equation (6.5) yields $\epsilon(A)<X$.  Also
\[
 2R(h_0-A)=\eta+D+(1-2R)A>0;
\]
when $R>1/2$, use $A<P=\eta E$ for the last sign.

Suppose, contrary to (6.14), that
\[
 A_0\le1-h_0.
 \tag{6.17}
\]
The companion edge is V-gap-free.  Thus
\[
 A_0+B_5\ge1,
\]
and (6.17) gives $B_5\ge h_0$.  Nonsupercriticality of $T_5$ and coverage
of the next center-free edge imply
\[
 \boxed{B_4\ge B_5\ge h_0.}
 \tag{6.18}
\]

The assumed transverse radial reach at $T_4$ is at least $1-A$.  Apply the exact local
catalogue to the reflected boundary pair $(B_4,A_4)$.  If the local point is
not on the selected $Q_+$ component, the constant, $Q_-$, and linear
components together with (6.16) give
\[
 1-A_4\ge1-\epsilon(A)>1-X.
\]
The center-free edges then carry this lower bound backward to
$B_1>1-X$, contradicting $A_1\ge X$ and
$A_1+B_1\le1$.

Hence the $T_4$ state is selected $Q_+$.  Lemma
\ref{lem:new-selected-chords} and (6.18) give
\[
 B_3>L_1,
 \qquad
 L_1=(2-4A)h_0-(1-4A)A.
 \tag{6.19}
\]
The selected-state condition also forces
\[
 \boxed{X>\frac12.}
 \tag{6.20}
\]
To see this, selectedness gives
\[
 h_0\le\epsilon(A)<\frac{2A}{1-2A}.
\]
If $X\le1/2$, the lower CE1 inequality gives
\[
 2Rh_0=\eta+A+D\ge w(R+A).
\]
Consequently
\[
 f_R(A)<0,
 \qquad
 f_R(z)=w(R+z)(1-2z)-4Rz.
\]
The quadratic is strictly concave and $f_R(0)>0$.  If $R\le1/2$, then
\[
 f_R(P)\ge\frac{1-E}{2}
 (6E^3-2E^2-5E+2)>0.
\]
If $R>1/2$, then $A<A_*=w/2-\eta$ and
\[
 f_R(A_*)=\frac{1-E}{2}(E+11R-3ER-5)>0.
\]
Concavity contradicts $f_R(A)<0$ in both cases, proving (6.20).  Thus
\[
 m<\frac w2<\frac14.
 \tag{6.21}
\]

If $B_3\ge1/2$, then $B_3>1-X$ and the preceding immediate contradiction
applies.  Assume $B_3<1/2$.  At $T_3$ the reflected incoming reach is
$B_3>h_0>m$, while its assumed transverse radial reach is at least $1-m$.  If this state is
not selected $Q_+$, the exact local catalogue gives
\[
 B_2\ge1-B_3>\frac12>1-X,
\]
and the remaining center-free edge again contradicts $T_1$.

Thus both $T_4$ and $T_3$ are selected $Q_+$.  A second use of Lemma
\ref{lem:new-selected-chords} gives
\[
 B_2>(2-5m)B_3-(1-5m)m.
\]
Since $2-5m>0$, (6.19) yields
\[
 \boxed{
 B_2>L_2,
 \qquad
 L_2=(2-5m)L_1-(1-5m)m.
 }
 \tag{6.22}
\]

We now prove the scalar estimate
\[
 \boxed{
 D<\frac1{10},
 \qquad
 L_2>2D+3D^2>\epsilon(D),
 \qquad
 \epsilon(D)<X.
 }
 \tag{6.23}
\]
The $T_4$ selected-state inequality and (6.6) give
\[
 D\le D_h:=A(4R-1)+10RA^2-\eta.
 \tag{6.24}
\]
If $D\ge1/10$, then $A+wD<P$ gives
\[
 A<\overline A=\frac{w(5R-1)}{10}.
\]
The right side of (6.24) is increasing in $A$, while $\eta>Rw/2$, so
\[
 D<U(R):=\overline A(4R-1)+10R\overline A^2-\frac{Rw}{2}.
\]
Direct expansion gives
\[
 \frac1{10}-U(R)=-\frac R{10}P_4(R),
\]
where
\[
 P_4(R)=25R^4-60R^3+26R^2+22R-14.
\]
For $y=2R-1\in(0,1)$,
\[
 16P_4(R)
 =25y^4-20y^3-106y^2+124y-39
 \le-101y^2+124y-39<0.
\]
Hence $U(R)<1/10$, a contradiction.  This proves $D<1/10$.

Put
\[
 \Psi(D)=L_2-2D-3D^2,
 \qquad
 J(D)=L_2-\frac{23}{10}D.
\]
Since $D<1/10$,
\[
 \Psi(D)-J(D)=3D\left(\frac1{10}-D\right)>0.
\]
For fixed $R,A$, differentiation of (6.22) gives
\[
 J'(D)=\frac{S(R,A)}{10R^2},
\]
where
\[
 S(R,A)=100A^2-(40R+50)A+R(20-23R).
\]
The nonempty interval $P-RA\le D\le D_h$ implies
\[
 A>\frac{4Rw}{25R-4}=:A_*.
\]
On the relevant interval $\partial S/\partial A<-45$, and
\[
 S(R,A)<S(R,A_*)
 =-\frac{Rq_3(R)}{(25R-4)^2},
\]
where
\[
 q_3(R)=8775R^3-14260R^2+7928R-1120.
\]
For $y=2R-1$,
\[
 8q_3(R)=8775y^3-2195y^2+997y+3007>0.
\]
Thus $J$ decreases with $D$, so $J(D)\ge J(D_h)$.

At $D=D_h$, one has $h_0=2A+5A^2$.  If $L_{2,h}$ is the corresponding value,
then
\[
 L_{2,h}-A(1+4R)=\frac{A}{R^2}Q(R,A),
\]
where
\begin{align*}
Q(R,A)={}&100A^3R-40A^2R^2-30A^2R+12AR^2-15AR+5A\\
&-4R^3+5R^2-R.
\end{align*}
This polynomial is concave in $A$ on $0\le A\le Rw/2$.  Its endpoint values
are
\[
 Q(R,0)=Rw(4R-1)>0
\]
and
\[
 Q\left(R,\frac{Rw}{2}\right)
 =\frac{Rw}{2}
 (25R^5-30R^4+20R^3-3R^2-7R+3)>0.
\]
The last positivity follows after putting $y=2R-1$:
\[
32p(R)=25y^5+65y^4+90y^3+106y^2-35y+5>0.
\]
Hence $L_{2,h}>A(1+4R)$.

Finally $A<P=\eta E$, $A<Rw/2$, and $1/E>1+Rw/2$ give
\[
 \frac{D_h}{A}
 <C_*:=4R-2+5R^2w-\frac{Rw}{2}.
\]
Moreover
\[
 1+4R-\frac{23}{10}C_*
 =\frac{230R^3-253R^2-81R+112}{20}>0.
\]
Thus $J(D_h)>0$, so
\[
 L_2>2D+3D^2.
\]
Equation (6.6) gives $\epsilon(D)<2D+3D^2$, and the argument proving
$\epsilon(A)<X$ also gives $\epsilon(D)<X$.  This proves (6.23).

The assumed transverse radial reach at $T_2$ is at least $1-D$.  Since
$B_2>\epsilon(D)$, Lemma~\ref{lem:new-direct-threshold}, applied after
reflection, gives
\[
 A_2\le\epsilon(D).
\]
Coverage of $e_{1,2}$ therefore forces
\[
 B_1\ge1-A_2\ge1-\epsilon(D)>1-X,
\]
contradicting $A_1\ge X$ and $A_1+B_1\le1$.  Hence (6.17) is impossible.
\end{proof}

\subsection{T3-like rescuer inequalities}
\label{subsec:calc-t3-rescuer}

With the notation of Proposition~\ref{prop:readable-one-t3-one-gap}, put
\[
 x=\frac{aD}{R_0},\qquad \theta=\frac{D-1}{R_0}.
\]
The midpoint constraints are
\[
 \frac{1-4\theta+\theta^2}{2(1-2\theta)}
 \le x\le\frac12-\theta,
 \qquad 0\le\theta\le2-\sqrt3.
\]
For $s(z)=z(2-z)/(1+z)$, the equation
$s(1-M_c^{\rm sup})=c=x+\theta$ shows that it is enough to prove
$s(x)\le x+\theta$.  After multiplying by $1+x$, this is
\[
 Q_\theta(x)=2x^2+(\theta-1)x+\theta\ge0.
\]
For $0\le\theta\le1/5$, substitution of the left endpoint gives
\[
 Q_\theta(x)\ge
 \frac{\theta(1-5\theta+11\theta^2-\theta^3)}
 {2(1-2\theta)^2}\ge0.
\]
For $1/5\le\theta\le2-\sqrt3$, the vertex lies in the feasible interval and
\[
 Q_\theta(x)\ge\frac{10\theta-1-\theta^2}{8}>0.
\]
Thus $x\le1-M_c^{\rm sup}$.  Since $D\ge R_0$,
\[
 a\le x,
 \qquad
 \frac{a}{a+1-u}=\frac{aD}{R_0}=x,
\]
which proves the two inequalities used in the body.

\subsection{Detailed nonzero-gap terminal calculations}
\label{subsec:new-detailed-direct-certificates}

This subsection supplies the calculations suppressed in the preceding case
proofs.  They are written as inequalities for actual reaches and radial
witnesses.  There is no formal chain notation.

\subsubsection{The disk--gap support calculation}
\label{subsec:calc-disk-gap}

Let $c=c_{\max}(p,q)$, $m=\min\{p,q\}$, and
$d=1-c$.  The regular radial hexagon has inradius $hd$.  Choose the farther
endpoint $G$ of the complementary gap.  Its norm is
\[
 \rho=\sqrt{1-m+m^2}.
\]
For a unit normal whose angle from $G$ is $\theta$, the three point supports
are
\[
 \rho\cos\theta,
 \quad
 \rho\cos(\theta+2\pi/3),
 \quad
 \rho\cos(\theta+4\pi/3).
\]
The disk replaces every one of these by at least $hd$.  A minimizing
orientation is either a disk-only orientation or a tie orientation.  At a tie
one projection equals $hd$.  The other active projection is
\[
 \frac{-hd+\sqrt{3(\rho^2-h^2d^2)}}2.
\]
Thus the support sum is at least $h$ exactly when
\[
 3d(1-d)\ge m(1-m).
 \tag{6.51}
\]
Since $d=1-c$, this is (6.29).

We now derive the local inequality used in the high-$c$ part.  Work in a
support orientation with angular parameter $t\in[0,\pi/3]$ and put
\[
 x=\cos\left|t-\frac\pi6\right|.
\]
The support inequalities for the local kite
$\conv\{0,pu,qv,c(u+v)\}$ imply
\[
 cx\le h
 \tag{6.52}
\]
and, after taking the smaller boundary demand $m$,
\[
 mx+\frac c2
 \left(x+\sqrt3\sqrt{1-x^2}\right)\le h.
 \tag{6.53}
\]
When $c>h$, the interval is $h\le x\le h/c$.  The left side of (6.53) is
concave.  If $c>1-m$, its value at $x=h$ is $h(m+c)>h$, so feasibility forces
its value at the other endpoint to be at most $h$:
\[
 \frac mc+\frac12+\sqrt{c^2-\frac34}\le1.
\]
Hence
\[
 m\le c\left(\frac12-\sqrt{c^2-\frac34}\right).
 \tag{6.54}
\]
Substitution of the right side into (6.51) gives
\[
\begin{aligned}
&3c(1-c)-
 c\left(\frac12-\sqrt{c^2-\frac34}\right)
 \left[1-c\left(\frac12-\sqrt{c^2-\frac34}\right)\right]\\
&\quad=
\frac{c(1-c)}2
\left(5-2c-2c^2+2\sqrt{c^2-\frac34}\right)\ge0.
\end{aligned}
\]
This proves the complementary-gap inequality directly from the four local
support conditions.

\subsubsection{The CE2 one-gap threshold inequalities}
\label{subsec:calc-ce2-one-gap-threshold}

We prove the two inequalities in (6.39).  Put
\[
 X=R-\delta,
 \qquad
 Q=\frac{\eta+\alpha+\delta}{2R}.
\]
The right-trace inequality gives
\[
 WX=RW-W\delta=\eta+P-W\delta>\eta+\alpha.
 \tag{6.55}
\]
The selected low root $\epsilon(\alpha)$ is the smaller root of
\[
 z^2-(1-\alpha)z+(1-\alpha)^2-(1-\alpha)^4=0.
 \tag{6.56}
\]
It is enough to prove that $X$ lies above that root.  Since
$0<X<R$ and the quadratic opens upward, substitute the lower estimate in
(6.55).  Equivalently consider
\[
 \pi(q)=(\eta+q)(1-2q)-2qW.
\]
This is concave.  At the endpoints of the admissible interval $0\le q\le P$,
\[
 \pi(0)=\eta>0,
\]
while the identities $RW=\eta+P$ and $P=E\eta$ give
\[
 \pi(P)=\eta(\eta+2ER^2)>0.
\]
Thus $\pi(\alpha)>0$, which is precisely
\[
 \epsilon(\alpha)<X.
 \tag{6.57}
\]

For the second inequality put $T=\alpha+\delta$.  Multiply
\[
 \alpha+W\delta<P,
 \qquad
 R\alpha+\delta<P
\]
by $W$ and $R$, respectively, and add.  Since
$W+R^2=W^2+R=E^2$,
\[
 E^2T<P=E\eta,
 \qquad
 T<\frac\eta E.
 \tag{6.58}
\]
Let $d=\min\{\alpha,\delta\}$.  Then $d\le T/2$.  By (6.5),
\[
 \min\{\epsilon(\alpha),\epsilon(\delta)\}
 <\frac{2d}{1-2d}
 \le\frac{T}{1-T}.
\]
It remains to prove $T/(1-T)<Q$.  This is equivalent to
\[
 2RT<(\eta+T)(1-T).
\]
Define
\[
 \chi(T)=(\eta+T)(1-T)-2RT.
\]
The function is concave.  At $T=0$, $\chi(0)=\eta>0$.  At
$T=\eta/E$, direct simplification gives
\[
 \chi\left(\frac\eta E\right)
 =\frac{\eta}{E^2}(E-R)^2>0.
\]
By (6.58), $T$ lies strictly between these endpoints.  Hence
\[
 \boxed{
 \min\{\epsilon(\alpha),\epsilon(\delta)\}<Q.
 }
 \tag{6.59}
\]
This proves the CE2 dichotomy without a composed reach map.

\subsubsection{The CE2 short-ray root separation}
\label{subsec:calc-short-ray}

We give the remaining details of Theorem~\ref{thm:new-ce2-short-ray}.  The
identities
\[
 RW=\eta+P,
 \qquad P=E\eta
\]
and the CE2 inequalities imply
\[
 Wq=W(R-\delta)>\eta+\alpha,
\]
\[
 Rp=R(W-\alpha)>\eta+\delta.
\]
Thus
\[
 p,q>\eta+e.
 \tag{6.60}
\]
Let $M=\max\{R,W\}$.  Since $\alpha,\delta\ge e$,
\[
 e<\frac{P}{1+M}.
\]
Moreover
\[
 (1+M)^2-3E^2=(2M-1)(2-M)\ge0,
\]
so
\[
 \eta>\sqrt3e.
 \tag{6.61}
\]
The universal bound
\[
 P\le\frac{2\sqrt3-3}{4},
 \qquad 1+M\ge\frac32
\]
gives
\[
 e<\frac{2\sqrt3-3}{6}=e_*.
\]

For $c=1-e$, the roots of
$f_c(t)=c^4-c^2+ct-t^2$ are $t_\pm$ from the proof above.  At
$t_0=(1+\sqrt3)e$,
\[
 f_c(t_0)=eB(e),
\]
with
\[
 B(e)=e^3-4e^2-3\sqrt3e+\sqrt3-1.
\]
Since
\[
 B'(e)=3e^2-8e-3\sqrt3<0
\]
on $[0,e_*]$, one has
\[
 B(e)\ge B(e_*)=\frac{302\sqrt3-501}{72}>0.
\]
Therefore $t_-<t_0<p,q$.  Also
\[
 p+q=1-\alpha-\delta\le1-2e=c-e.
\]
If, say, $q\le p$, then $q>t_-$ and
\[
 p<c-e-t_-<c-t_-=t_+.
\]
The same holds with $p,q$ interchanged.  Hence
\[
 t_-<p,q<t_+,
\]
which proves $f_c(p),f_c(q)>0$.  In the exact local finite-caliper formula at
radial demand $c$, the four side candidates are
\[
 p+c,
 \quad q+c,
 \quad\frac{c^2}{\sqrt{p^2-pc+c^2}},
 \quad\frac{c^2}{\sqrt{q^2-qc+c^2}}.
\]
The first two exceed one because $p,q>e$, and the last two exceed one exactly
because $f_c(p),f_c(q)>0$.  This completes the direct short-ray proof.

\subsubsection{Detailed adjacent Vd residual estimate}
\label{subsec:calc-adjacent-vd}

We justify (6.48).  In the adjacent placement define the residuals
$\rho_R,\rho_L$ as in Case 1 of Proposition~\ref{prop:new-one-vd-assembly}.
The signed CE2 domain first gives
\[
 \alpha+\delta<\frac16\min\{R,W\}.
 \tag{6.62}
\]
The positive residual $\rho_L$ has one of two forms.

If
\[
 \rho_L=W-\alpha,
\]
then
\[
 4\delta+\alpha
 \le4(\alpha+\delta)
 <\frac{2W}{3}<W,
\]
so $4\delta<\rho_L$.

Assume instead
\[
 \rho_L=1-A_0.
\]
The other residual cannot equal $1-B_0$, since (6.47) would imply
$A_0+B_0>3/2$, contrary to the endpoint-distance bound
$A_0+B_0\le2/\sqrt3$.  Hence
\[
 \rho_R=1-(W+\delta),
 \qquad
 B_0\ge y:=\frac{k}{R}.
\]
The two points of $T_0$ with reaches $1-\rho_L$ and $y$ have distance at most
one.  The diameter-transfer inequality gives
\[
 \rho_L>\frac y2>\frac{\eta}{2R}.
 \tag{6.63}
\]
Since $\rho_R+\rho_L<1/2$,
\[
 \delta>R-\frac12+\frac{\eta}{2R}=:J(R).
 \tag{6.64}
\]
Using $\eta/R=W/(1+E)$, exact differentiation gives
\[
4E(1+E)^2J'(R)
=2E(1+E)(1+2E)-W(2R-1)>0.
\]
Thus $J$ is strictly increasing.  Since $J(3/8)=1/24$ and (6.62) gives
$\delta<1/24$, equation (6.64) forces $R<3/8$.  Then
\[
 (2-3R)^2-E^2=(3-8R)(1-R)>0,
\]
so $E<2-3R$ and
\[
 \frac{\eta}{R}=\frac{W}{1+E}>\frac13.
\]
Together with (6.63),
\[
 \rho_L>\frac16>4\delta.
\]
This proves (6.48) in both residual cases.

\subsubsection{Detailed nonadjacent Vd separation}
\label{subsec:calc-nonadjacent-vd}

Let $x=k/W$ and $y=k/R$ be the two center near endpoints.  Let $B,U$ be the
farthest already-covered extents on the corresponding edges.  We prove
\[
 B\le M_0(x),
 \qquad
 U\le M_0(y).
 \tag{6.65}
\]
By symmetry it is enough to treat $B$.  The residual component formula says
that $B$ is either $W+\delta$ or the supercritical endpoint $B_0$.
If $B=W+\delta$, the center contains boundary points with reaches $x,B$, so
the diameter-one inequality gives $B\le M_0(x)$.  If $B=B_0$ and
$A_0<x$, then the opposite covered extent is $U=A_0$; equation (6.47) would
again force $A_0+B_0>3/2$.  Hence $A_0\ge x$, and monotonicity of $M_0$
gives
\[
 B_0\le M_0(A_0)\le M_0(x).
\]
This proves (6.65).

The elementary identity
\[
 1-M_0(z)>\frac z2
 \qquad(z>0)
\]
then gives
\[
 \rho_R>\frac x2,
 \qquad
 \rho_L>\frac y2.
\]
The common center endpoint-slack bound is
\[
 \alpha+\delta<\frac12\min\{x,y\}.
\]
Thus (6.49) follows.  The exact Vd corner form gives the strict own-radial
bound used in Case 2.

\subsubsection{Detailed Vd1 rescue inequalities}
\label{subsec:calc-vd1-rescuer}

Use the Vd1 corner parameter $t\ge1$, let
$d=\sqrt{t^2+t+1}$, and write the supported interval as $[c,u]$.  The endpoint
relations are
\[
 c=\frac{t(1-b)}{t+1},
 \qquad
 u=\frac{d-a-tb-1}{t}.
\]
The midpoint condition gives
\[
 a+tb\le d-1-\frac t2.
\]
Eliminating $tb=t-c(t+1)$ gives
\[
 a\le F(t,c):=d-1-\frac{3t}{2}+c(t+1).
\]
For fixed $c\le1/2$,
\[
 \frac{\partial F}{\partial t}
 =\frac{2t+1}{2\sqrt{t^2+t+1}}-\frac32+c<0,
\]
so
\[
 a\le F(1,c)=\sqrt3-\frac52+2c=:L(c).
\]
The inequality $2L(c)\le1-M_c^{\rm sup}$ is equivalent, after squaring, to
\[
 Q(c)=20c^2+(18\sqrt3-52)c+47-24\sqrt3\ge0.
\]
The polynomial decreases on $[0,1/2]$ and
$Q(1/2)=26-15\sqrt3>0$.  Hence
\[
 a\le1-M_c^{\rm sup}.
\]

Put $\varepsilon=1-u$.  Direct substitution gives
\[
 \varepsilon=\varepsilon_0+\frac at,
 \qquad
 \varepsilon_0=
 \frac{2t+1-c(t+1)-d}{t}>0.
\]
The function
\[
 z\longmapsto\frac{z}{z+\varepsilon_0+z/t}
\]
is increasing.  Since
\[
 \varepsilon_0+\frac{F(t,c)}t=\frac12,
\]
we obtain
\[
 \frac{a}{a+1-u}
 \le\frac{F(t,c)}{F(t,c)+1/2}
 \le2L(c)
 \le1-M_c^{\rm sup}.
\]
This proves the two direct Vd1 rescue inequalities in (6.50).

\subsection{Casewise witness audit}
\label{subsec:new-casewise-witness-details}

We now record the finite set used in each combinatorial row, including the
open-endpoint issues.  This makes the Strategy~3 ownership independent of the
former reach-propagation chapter.

\subsubsection{One gap and no supercritical V role}

Normalize the only gap to
\[
 J=[B_0,1-A_1]=[\ell,r]\subset e_{0,1}.
\]
On every other edge choose an open-overlap point with coordinate $x_i$.
Because $A_i+B_i\le1$, one has
\[
 x_1\le r,
 \qquad
 x_i\le x_{i-1}\quad(2\le i\le5),
 \qquad
 x_5\ge\ell.
 \tag{6.72}
\]
The endpoint inequalities remain valid when a gap or overlap degenerates to a
singleton; the strict open choices are obtained by an arbitrarily small
perturbation and the final enclosure number is continuous.

Put $p=1-r$ and $q=\ell$.  For $T_1$, the two selected lower demands are
$p$ and $x_1\ge q$.  For $2\le i\le5$, they are
$1-x_{i-1}\ge p$ and $x_i\ge q$.  At $T_0$, they are
$1-x_5\ge p$ and $q$.  Thus every local role dominates $(p,q)$.

Let $c_*=c_{\max}(p,q)$.  The own role cannot contain
$(1-c_*)V_i$ by the definition of $c_{\max}$.  If an adjacent role is Vd0,
it has no positive interval on $r_i$.  If it is T3-like, its unique adjacent
capacity is $C_+$ or $C_-$ and is at most $c_*$ by
Lemma~\ref{lem:new-common-pair-domination}.  The other roles are excluded by
diameter.  Hence
\[
 (1-c_*)V_i\in U_C
 \qquad(0\le i\le5).
 \tag{6.73}
\]
The six points form a regular hexagon of circumradius $1-c_*$.  The convex
center role therefore contains its inball
\[
 \mathcal D_{h(1-c_*)}.
\]
The gap segment is exactly $J(p,q)$.  This proves the hypotheses of
Theorem~\ref{thm:new-complementary-gap} without any abstract intermediate-set
argument.

\subsubsection{The T3-like neighboring capacity}

We verify directly that the T3-like supported endpoint is the radical branch
of Proposition~\ref{prop:new-neighbor-ray-formula}.  Use the translated
T3-like normal form with boundary reaches $a,b$ and shape parameter $D$.
The support-side relation is
\[
 R_0^2-DR_0+D^2=1,
\]
and the far endpoint on the forward neighboring ray is
\[
 u=1+a-\frac{R_0}{D}.
\]
The boundary reaches satisfy
\[
 a+b=\frac1D.
\]
Solving the quadratic for $R_0/D$ gives
\[
 \frac{R_0}{D}
 =\frac{1+\sqrt{4(a+b)^2-3}}2.
\]
Therefore
\[
 \boxed{
 u=a+\frac12-\sqrt{(a+b)^2-\frac34},
 }
 \tag{6.74}
\]
which is exactly the final radical branch of $C_+(a,b)$.  Thus no additional
T3-specific radial function occurs in the one-gap proof.

A T3-like role is nonsupercritical.  Hence its boundary sum is at most one,
and (6.74) is bounded by
\[
1-\min\{a,b\}\le c_{\max}(p,q)
\]
whenever $(a,b)$ dominates the common pair $(p,q)$.  This is the explicit
reason why one or two T3-like roles do not shrink the regular radial witness
hexagon in Proposition~\ref{prop:new-nplus-zero-gap-closures}.

\subsubsection{Two gaps}

For CE2, write the two center intervals as
\[
 I_L=\left[\frac{k}{W},R+\alpha\right],
 \qquad
 I_R=\left[\frac{k}{R},W+\delta\right].
\]
If both contain gaps, then the incident endpoint roles give
\[
 A_1\ge R-\delta=q,
 \qquad
 B_5\ge W-\alpha=p.
\]
Choose handoffs $x_1,\ldots,x_4$ on the four intervening edges.  The five
nonsupercritical roles give
\[
 p<x_4<x_3<x_2<x_1<1-q.
 \tag{6.75}
\]
Every one of $T_1,\ldots,T_5$ therefore dominates $(p,q)$.  In particular,
the two transverse witnesses
\[
 D_2=(1-c_{\max}(p,q))V_2,
 \qquad
 D_4=(1-c_{\max}(p,q))V_4
\]
are excluded from all V roles.  The center exits at distances
$\delta$ and $\alpha$ on these two rays.  Theorem
\ref{thm:new-ce2-short-ray} states that one witness lies strictly beyond its
exit.  This argument does not inspect $A_0+B_0$, so it simultaneously handles
$N_+=0$ and $N_+=1$.

\subsubsection{One gap, one supercritical Vd0 role}

In the all-Vd0 branch, define the six actual radial endpoints
\[
 P_i=(1-C_i)V_i.
\]
Because $U_i$ is open, the endpoint does not belong to $U_i$.  The adjacent
roles have no positive support on $r_i$, and nonlocal roles are separated by
distance.  Hence $P_i\in U_C$.

Suppose a second open unit triangle $U'$ contains
\[
 O,M_0,X(\ell),X(r),P_0,\ldots,P_5.
\]
The fixed six V roles together with $U'$ cover every radial arm: $U'$ contains
$[O,P_i]$, and $U_i$ contains a relatively open interval from $V_i$ past every
point short of $P_i$.  They also cover the perimeter, since the unique missing
boundary component is the gap and $U'$ contains its two endpoints and hence
its full segment.  Thus all direct midpoint and signed-center conclusions
apply to $U'$.

If $d_i'$ is the exit of $U'$ on $r_i$, containment of $P_i$ says
\[
 d_i'\ge1-C_i,
 \qquad
 C_i\ge1-d_i'.
 \tag{6.76}
\]
This is the only connection between the finite set and the local reach
calculation.  In CE2, (6.76) supplies radial reaches $1-\alpha$ at $T_4$ and
$1-\delta$ at $T_2$; the direct threshold alternatives of the proof of
Proposition~\ref{prop:new-nplus-one-all-vd0} contradict the terminal bound at
$T_0$.  In CE1, (6.76) supplies the reaches used by Proposition
\ref{prop:new-ce1-direct-certificate}.  No statement about a residual polygon
is substituted for these actual endpoint facts.

For a singleton gap, replace $r=\ell$ by $r_n>\ell$ with $r_n\downarrow\ell$.
The diameter inequalities and the unique-supercritical pattern persist for
large $n$.  The witness sets converge in the Hausdorff metric, and
$\Lambda$ is continuous by the support formula.  Hence the positive-length
proof passes to the singleton limit.

\subsubsection{The T3-like rescuer endpoint}

In Proposition~\ref{prop:new-one-t3-terminal}, the finite point
$P_T=(1-u)V_1$ is forced into the center for a stronger reason than a generic
radial capacity.  The T3-like open trace begins at the vertex side at $c$ and
ends at $u$, so its O-side endpoint is not in the open role.  The adjacent
supercritical role misses $M_1$ and therefore, by convexity, misses the entire
O-side segment.  All remaining roles are Vd0.  Thus the segment from $O$ to
$P_T$ is owned by the center.

The two inequalities (6.44) then compare the far boundary tail with the one
free strict-supercritical envelope.  They do not propagate a numerical state
through the four middle roles.  Instead, adding the four actual edge-cover
inequalities gives one path budget:
\[
(A_2+B_2)+(A_3+B_3)+(A_4+B_4)+(A_5+B_5)
\ge4+h_T-B_1>4.
\]
This direct sum is the full global contradiction.

\subsubsection{The Vd radial points}

In the adjacent Vd placement, define
\[
 q_2=1-\delta.
\]
The two roles capable of reaching the center interval on $r_2$ have exact
far endpoints bounded by
\[
 u_{1\to2}<q_2,
 \qquad
 C_2<q_2.
\]
Therefore any point of $r_2$ between the larger of these two endpoints and
$q_2$ is missed by all seven roles.  This is a literal one-point obstruction,
not a perimeter reformulation.

In the nonadjacent placement, the point
\[
 D_\tau=\min\{\rho_R,\rho_L\}V_\tau
\]
is beyond the Vd own-radial trace by (6.70), and beyond the center trace by
(6.49).  Its adjacent roles are Vd0 and its nonlocal roles are separated by
diameter.  This is again a literal uncovered point.

In the Vd1 rescue placement, the O-side endpoint of the supported interval
plays exactly the same role as $P_T$ in the T3-like case.  Only the local
normal-form proof of the two inequalities changes; the center-hiding and
four-role boundary sum are identical.

\subsubsection{The corrected replacement placement}

When neither distinguished role is based at $V_0$ and the exceptional role is
Vd1, the shared edge is center-free.  The corrected replacement constructs
one nonsupercritical Vd0 role in each of the two correct vertex charts.  It
preserves all six boundary and radial skeleton traces needed for coverage but
does not claim to preserve the number of gaps.  After replacement, recompute
$N'_{\rm gap}$.

If $N'_{\rm gap}=0$, the six nonsupercritical Vd0 roles cover the perimeter,
contradicting the strict boundary-overlap sum.  If $N'_{\rm gap}=1$, use the
common radial disk and complementary gap.  If $N'_{\rm gap}=2$, use the two
transverse short-ray witnesses.  Hence the replacement terminal has no return
to a reach-composition argument.

\subsection{Selector and singleton-gap audit}
\subsubsection{Branchwise proof of common-pair domination}

We expand the neighboring-capacity comparison in Lemma
\ref{lem:new-common-pair-domination}.  Assume first that $q\le p$.
The own-radial demand $1-q$ is admissible because the exact edge-normal
orientation has side
\[
 q+\max\{p,1-q\}=1.
\]
Thus
\[
 c_{\max}(p,q)\ge1-q.
 \tag{6.73a}
\]
For $C_+(p,q)$ there are three possibilities.

On the linear branch,
\[
 C_+(p,q)=1-q.
\]
On the plateau branch, $C_+=p(p)$, where the cubic root satisfies
\[
 f_p(1-q)
 =(1-q)^3-(p+2)(1-q)^2+2(p+1)(1-q)-1\ge0
\]
throughout the selected interval.  Since the cubic is increasing,
$p(p)\le1-q$.  On the radical branch, the capacity decreases from its
selected transition value and therefore is no larger than the plateau.  Hence
\[
 C_+(p,q)\le1-q.
\]
The reflected calculation gives the same bound for $C_-$ and treats the order
$p\le q$.  Combining with (6.73a) proves
\[
 C_\pm(p,q)\le c_{\max}(p,q).
\]

The same argument is valid for any actual pair $(a,b)$ with $a\ge p$ and
$b\ge q$, because the family of triangles containing the larger anchors is a
subfamily of the one containing the smaller anchors.  Thus the relevant
capacity is coordinatewise nonincreasing.

\subsubsection{Finite-caliper selector audit}

For a finite point set $K$, Proposition~\ref{prop:new-enclosure-gauge} reduces
$\Lambda(K)$ to finitely many orientations.  We spell out the reduction used
for every diagram and for the anisotropic one-gap set.  Let the convex hull
vertices be $P_0,\ldots,P_{r-1}$ in cyclic order.  On an angular interval on
which the three support points are fixed, write
\[
 H(\theta)=
 \sum_{j=0}^2
 \langle Q_j,n_{\theta+2\pi j/3}\rangle.
\]
Then
\[
 H''(\theta)=-H(\theta).
\]
The support sum is positive unless $K$ is a singleton at the origin.  Hence an
interior critical point is a strict maximum.  A minimum occurs when two
support points tie in one normal direction.  That normal is perpendicular to
a hull edge $P_iP_{i+1}$.  Therefore
\[
 \Lambda(K)=
 \min_{0\le i<r}
 \frac1h
 \sum_{j=0}^2
 h_K(\mathsf R^jn_i),
 \tag{6.73b}
\]
where $n_i$ is either unit normal to $P_iP_{i+1}$.

This finite-caliper identity has two consequences.  First, the enclosure
number is unchanged by replacing a witness set by its convex hull.  Second,
no numerical angular minimization is hidden in the proof: every terminal can
be checked at finitely many support contacts.  The complementary-gap theorem
uses this fact analytically; the zero-gap nine-point theorem uses it through
the support-arc certificate in the final appendix.

\subsubsection{Terminal upper bound in the one-gap exact-one case}

We give the elementary derivation of (6.38).  The active center trace begins
at $k/R$, so the unique supercritical role has forward boundary reach
\[
 B_0\ge\frac{k}{R}.
\]
The center exit on $r_0$ is $1-k$.  Since the finite set contains the actual
radial endpoint $P_0=(1-C_0)V_0$, containment by the candidate center gives
\[
 C_0\ge k.
\]
The closure of $T_0$ therefore contains three points based at $V_0$ with
selected reaches
\[
 A_0,
 \qquad \frac{k}{R},
 \qquad k.
\]
Discarding the radial point only weakens the diameter condition, and gives
\[
 A_0\le M_0\left(\frac{k}{R}\right).
\]
For $0<z<1$,
\[
 M_0(z)<1-\frac z2.
\]
Indeed, after moving terms and squaring, this is equivalent to
$3z^2>0$.  With $z=k/R$,
\[
 A_0<1-\frac{k}{2R}=1-Q.
\]
This is the terminal capacity contradicted separately by the direct CE1 and
CE2 reach arguments.

\subsubsection{Strictness at singleton gaps}

Suppose the V-gap on $e_{0,1}$ is the singleton $X(\ell)$.  Although its
length is zero, neither incident open V trace contains the point.  Boundary
locality therefore forces $X(\ell)\in U_C$.  Choose
$r_n>\ell$ decreasing to $\ell$, and reduce the incoming reach at $T_1$ from
$1-\ell$ to $1-r_n$.  Every local diameter inequality remains feasible, and
the strict supercritical pattern is unchanged for large $n$.  The non-singleton-gap
witness sets converge to the singleton witness set.  By Lemma
\ref{lem:new-lambda-continuity}, the inequality $\Lambda\ge1$ passes to the
limit.  Thus no proof in this section charges numerical gap length; every one
uses the actual missing point.

\subsubsection{Why a common disk is insufficient for $N_+=1$}

The one-gap $N_+=1$ finite set retains all six actual radial endpoints rather
than replacing them by one common radius.  This is necessary.  The boundary
sequence has two ascents when the gap jump is included: the V-gap jump and the
unique supercritical V triangle.  They need not join the global minimum to the global
maximum.  A common lower pair can therefore move several radial witnesses
toward $O$ and can produce a disk-plus-gap set with enclosure number below
one.  The anisotropic set $P_0,\ldots,P_5$ preserves the six actual radial
constraints, and equations (6.76), (6.38), and the direct CE1/CE2 arguments
use those constraints individually.  This is the precise point at which the
new proof differs from a relabeling of the former propagation calculation.
